% ---
% id: statistics/binomial-box/lem:envelope-dominance
% kind: lem
% title: "Binomial dominance under bounded conditional success probabilities"
% summary: >-
%   A count of successes whose conditional success probability in each round,
%   given the past, lies in [a, b] has an upper tail no heavier than the
%   binomial law at b and a lower tail no heavier than that at a. The box can
%   therefore be built from binomial tails without independent rounds.
% prominence: standard
% status: accepted
% uses: []
% mentions: [thm:completeness-bound, statistics/ass:hon-env]
% assumes: []
% refs: []
% provenance: newly-derived
% lean:
% notes: "It needs no assumption: it is a statement about any adapted sequence; A-HON-ENV is what makes the per-outcome indicators of an honest run such a sequence (thm:completeness-bound)."
% ---

This is the classical stochastic dominance of a sum of adapted Bernoulli
variables by a binomial law, derived here in full by a backward induction on
the conditional tail. Under Assumption~\ref{ass:hon-env} the indicators of
one outcome over the rounds of an honest run are such a sequence, with $a$
and $b$ the envelope's two values at that outcome (Theorem~\ref{thm:completeness-bound}).

\begin{lemma}\label{lem:envelope-dominance}
Let $\mathcal{F}_0\subseteq\mathcal{F}_1\subseteq\dots\subseteq\mathcal{F}_n$
be a filtration on a probability space, $n\ge1$, and let
$\chi_1,\dots,\chi_n$ be $\{0,1\}$-valued random variables with $\chi_j$
$\mathcal{F}_j$-measurable. Let $0\le a\le b\le1$ and suppose that, for every
$j\in\{1,\dots,n\}$, almost surely
\begin{equation}\label{eq:dominance-hyp}
  a\ \le\ P_j\ \coloneqq\ \Pr\big[\chi_j=1\ \big|\ \mathcal{F}_{j-1}\big]\ \le\ b .
\end{equation}
Put $S\coloneqq\sum_{j=1}^{n}\chi_j$, and for $m\in\{0,\dots,n\}$,
$r\in[0,1]$ and an integer $k$ write
$G^{r}_{m}(k)\coloneqq\Pr[\mathrm{B}^{r}_{m}\ge k]$, where
$\mathrm{B}^{r}_{m}$ is binomial with parameters $(m,r)$ ($\mathrm{B}^{r}_{0}=0$).
Then for every integer $k$
\begin{equation}\label{eq:dominance}
  \Pr[S\ge k]\ \le\ \Pr\big[\mathrm{B}^{b}_{n}\ge k\big],
  \qquad
  \Pr[S\le k]\ \le\ \Pr\big[\mathrm{B}^{a}_{n}\le k\big],
\end{equation}
and equivalently
\begin{equation}\label{eq:dominance-strict}
  \Pr[S>k]\ \le\ \Pr\big[\mathrm{B}^{b}_{n}>k\big],
  \qquad
  \Pr[S<k]\ \le\ \Pr\big[\mathrm{B}^{a}_{n}<k\big].
\end{equation}
\end{lemma}

\begin{proof}
\emph{Two facts about the binomial tail.} (1) For $m\ge0$ and every integer
$k$, $G^{r}_{m}(k-1)-G^{r}_{m}(k)=\Pr[\mathrm{B}^{r}_{m}=k-1]\ge0$. (2) The
\emph{Pascal recursion of the binomial tail}: for $m\ge0$ and every integer
$k$,
\begin{equation}\label{eq:pascal-tail}
  G^{r}_{m+1}(k)\ =\ r\,G^{r}_{m}(k-1)+(1-r)\,G^{r}_{m}(k).
\end{equation}
Indeed $\mathrm{B}^{r}_{m+1}$ has the law of $X+Y$ with $Y$ binomial
$(m,r)$ and $X$ Bernoulli$(r)$ independent of $Y$, and conditioning on $X$
gives $\Pr[X+Y\ge k]=r\Pr[Y\ge k-1]+(1-r)\Pr[Y\ge k]$.

\emph{The upper tail, in conditional form.} For $j\in\{0,\dots,n\}$ put
$S_{>j}\coloneqq\sum_{i=j+1}^{n}\chi_i$ (so $S_{>0}=S$ and $S_{>n}=0$). We show
by backward induction on $j$ that for every integer $k$, almost surely,
\begin{equation}\label{eq:dominance-cond}
  \Pr\big[S_{>j}\ge k\ \big|\ \mathcal{F}_j\big]\ \le\ G^{b}_{n-j}(k).
\end{equation}
\emph{Base $j=n$.} $S_{>n}=0$, so the left side is $\mathbf{1}\{k\le0\}
=G^{b}_{0}(k)$.
\emph{Step from $j$ to $j-1$}, for $j\in\{1,\dots,n\}$, with $m\coloneqq n-j$.
Since $S_{>j-1}=\chi_j+S_{>j}$ and $\chi_j\in\{0,1\}$,
\begin{equation*}
  \mathbf{1}\{S_{>j-1}\ge k\}
  =\chi_j\,\mathbf{1}\{S_{>j}\ge k-1\}+(1-\chi_j)\,\mathbf{1}\{S_{>j}\ge k\}.
\end{equation*}
Conditioning on $\mathcal{F}_j$, with $\chi_j$ $\mathcal{F}_j$-measurable and
bounded, and using the induction hypothesis at $k-1$ and at $k$ together with
$\chi_j\ge0$ and $1-\chi_j\ge0$, gives almost surely
\begin{equation*}
  \Pr\big[S_{>j-1}\ge k\,\big|\,\mathcal{F}_j\big]
  \le\chi_j\,G^{b}_{m}(k-1)+(1-\chi_j)\,G^{b}_{m}(k).
\end{equation*}
Conditioning further on $\mathcal{F}_{j-1}\subseteq\mathcal{F}_j$ (the tower
property, and monotonicity of conditional expectation) gives almost surely
\begin{equation*}
  \Pr\big[S_{>j-1}\ge k\,\big|\,\mathcal{F}_{j-1}\big]
  \le G^{b}_{m}(k)+P_j\,\big(G^{b}_{m}(k-1)-G^{b}_{m}(k)\big)
  \le G^{b}_{m}(k)+b\,\big(G^{b}_{m}(k-1)-G^{b}_{m}(k)\big)
  =G^{b}_{m+1}(k),
\end{equation*}
the second inequality by $P_j\le b$ from \eqref{eq:dominance-hyp} and fact
(1), the equality by \eqref{eq:pascal-tail}. As $m+1=n-(j-1)$, this is
\eqref{eq:dominance-cond} at $j-1$. At $j=0$, taking expectations gives
$\Pr[S\ge k]\le G^{b}_{n}(k)$, the first inequality of \eqref{eq:dominance}.
Only the upper bound $b$ was used.

\emph{The lower tail.} Put $\chi'_j\coloneqq1-\chi_j$. These are
$\{0,1\}$-valued and adapted, and
$\Pr[\chi'_j=1\,|\,\mathcal{F}_{j-1}]=1-P_j\le1-a$ almost surely by
\eqref{eq:dominance-hyp}. The upper-tail statement applied to them with the
bound $1-a$ gives, for every integer $k$,
$\Pr[S\le k]=\Pr[\,n-S\ge n-k\,]\le\Pr[\mathrm{B}^{1-a}_{n}\ge n-k]$. Since
$n-\mathrm{B}^{a}_{n}$ is binomial with parameters $(n,1-a)$ (it counts the
failures), the right side is $\Pr[n-\mathrm{B}^{a}_{n}\ge n-k]
=\Pr[\mathrm{B}^{a}_{n}\le k]$, the second inequality of
\eqref{eq:dominance}.

\emph{The strict forms.} $S$ and the binomial variables are integer-valued, so
$\{S>k\}=\{S\ge k+1\}$ and $\{S<k\}=\{S\le k-1\}$, and likewise for
$\mathrm{B}^{b}_{n}$ and $\mathrm{B}^{a}_{n}$; \eqref{eq:dominance} at $k+1$
and at $k-1$ is \eqref{eq:dominance-strict}, and conversely.
\end{proof}
